\documentclass[11pt]{article}

\usepackage{nmrdoc}
\usepackage{siunitx}
\usepackage{bm}

\sisetup{per-mode=symbol}
\DeclareSIUnit{\angstrom}{\text{\AA}}
\DeclareSIUnit{\bohr}{a_0}
\DeclareSIUnit{\hartree}{Ha}
\DeclareSIUnit{\elementarycharge}{e}

\begin{document}
\NmrTitle{The EEQ Charge-Equilibration Calculator}

\section{What the calculator computes}

The EEQ calculator assigns a geometry-dependent partial charge \(q_i\) to each
atom \(i\) in a protein conformation. It uses the extended electronegativity
equalization model: atoms exchange continuous charge until the marginal energy
cost of putting a little more charge on any atom is the same, while the total
charge of the conformation is held fixed. In this implementation the default
total charge is \(Q=\SI{0}{\elementarycharge}\), so the stored charge vector is
neutral unless the later per-atom clamp is triggered.

These charges bear on shielding through electric-field effects. A distribution
of charge produces an electric field \(\bm E\) and an electric-field gradient
\(\bm V\) at a nucleus. In Buckingham's expansion, the electronic shielding can
change with the local electric field, and the anisotropic part can couple to the
rank-2 field-gradient tensor. That is the physical route from a charge model to
an NMR shielding descriptor.

The code for \texttt{EeqResult} stops before the field calculation. It writes
\texttt{eeq\_charges.npy}, an \(N\)-component vector of charges in elementary
charge units, and \texttt{eeq\_cn.npy}, the coordination numbers used while
building those charges. It does not emit a shielding tensor, an electric field,
an electric-field gradient, or a polarizability. The output is therefore a
geometry-derived charge descriptor, not a chemical shift and not a shielding in
ppm.

\section{Where shielding enters}

For a nucleus in an applied magnetic field \(\bm B_0\), the local field is
usually written
\begin{equation}
  \bm B_{\mathrm{nuc}} = (\bm I-\bm\sigma)\bm B_0 .
  \label{eq:shielding}
\end{equation}
\(\bm I\) is the \(3\times3\) identity matrix, and \(\bm\sigma\) is the nuclear
shielding tensor. In physical NMR \(\bm\sigma\) is dimensionless: its entry
\(\sigma_{ab}\) says how a field applied along Cartesian direction \(b\) changes
the field component felt by the nucleus along direction \(a\). The tensor form
is needed because a molecule has orientation. Rotating the same local structure
can change both the magnitude and the direction of the secondary magnetic field
created by the electrons.

The quantum-mechanical origin is electronic response. The magnetic field
perturbs the electronic state; the perturbed electrons create their own field at
the nucleus; Eq.~\eqref{eq:shielding} records the linear coefficient between
the applied field and the field that reaches the nucleus. Electric fields enter
because the electron cloud is also shaped by the electrostatic environment. A
nearby charge distribution can pull electron density toward one side of a bond
or lone pair, changing the magnetic response of that electronic state.

For point charges \(q_j\), the vacuum field and field gradient at a target atom
\(i\) would have the Coulomb forms
\begin{align}
  E_a(i) &=
    k_e \sum_{j\ne i} q_j
    \frac{(r_i-r_j)_a}{|\bm r_i-\bm r_j|^3},
    \label{eq:field}\\
  V_{ab}(i) &=
    k_e \sum_{j\ne i} q_j
    \left[
      \frac{3(r_i-r_j)_a(r_i-r_j)_b}{|\bm r_i-\bm r_j|^5}
      - \frac{\delta_{ab}}{|\bm r_i-\bm r_j|^3}
    \right].
    \label{eq:efg}
\end{align}
Here \(\bm r_i\) is the target position, \(\bm r_j\) is a source-charge
position, \(a,b\in\{x,y,z\}\), \(\delta_{ab}\) is one when \(a=b\) and zero
otherwise, and \(k_e=\SI{14.3996}{\volt\angstrom}\) converts the raw
\(\si{\elementarycharge\per\angstrom\squared}\) and
\(\si{\elementarycharge\per\angstrom\cubed}\) sums to
\(\si{\volt\per\angstrom}\) and
\(\si{\volt\per\angstrom\squared}\). Equations~\eqref{eq:field} and
\eqref{eq:efg} explain why a charge model matters for shielding, but they are
not executed by \texttt{EeqResult.cpp}.

\section{The tensor split used by the pipeline}

When a calculator does produce a \(3\times3\) tensor \(\bm X\), the shared
\texttt{SphericalTensor} code decomposes it into scalar, antisymmetric, and
symmetric-traceless pieces. The scalar part is
\begin{equation}
  T_0 = \frac{1}{3}\operatorname{tr}\bm X .
  \label{eq:t0}
\end{equation}
This is the isotropic value: it multiplies the identity matrix and is the part
that survives orientational tumbling in solution.

The antisymmetric part is stored as a three-component pseudovector,
\begin{equation}
  \bm T_1 =
  \frac{1}{2}
  \begin{pmatrix}
    X_{yz}-X_{zy}\\
    X_{zx}-X_{xz}\\
    X_{xy}-X_{yx}
  \end{pmatrix}.
  \label{eq:t1}
\end{equation}
This is a compact representation of the part of \(\bm X\) that changes sign
under transpose. It is usually not observed directly in standard isotropic
solution NMR.

The remaining part is the symmetric traceless matrix
\begin{equation}
  \bm S =
  \frac{\bm X+\bm X^{\mathsf T}}{2}
  - T_0\bm I,\qquad \operatorname{tr}\bm S=0 .
  \label{eq:s}
\end{equation}
\(\bm S\) has five independent entries. The code packs them as
\begin{equation}
  \bm T_2 =
  \begin{pmatrix}
    \sqrt{2}\,S_{xy}\\
    \sqrt{2}\,S_{yz}\\
    \sqrt{3/2}\,S_{zz}\\
    \sqrt{2}\,S_{xz}\\
    (S_{xx}-S_{yy})/\sqrt{2}
  \end{pmatrix},
  \label{eq:t2}
\end{equation}
in the order \(m=-2,-1,0,+1,+2\). These factors are the normalization in
\texttt{Types.cpp}; they preserve the Frobenius norm, so
\(\sum_m T_{2m}^2=\sum_{ab}S_{ab}^2\). The packed nine-value layout is
\[
  [T_0,\;T_1[0],T_1[1],T_1[2],\;T_2[0],T_2[1],T_2[2],T_2[3],T_2[4]] .
\]

EEQ does not populate this tensor layout. A partial charge is a scalar attached
to an atom. It can be treated as a rotational rank-0 feature in a later model,
but it is not the \(T_0\) part of a shielding tensor. The anisotropic rank-2
argument enters when charges are turned into a field-gradient tensor such as
Eq.~\eqref{eq:efg}; the present EEQ source file does not do that step.

\section{The method implemented}

The implementation first reads atom elements and coordinates from the
conformation. Coordinates are converted from \(\si{\angstrom}\) to \(\si{\bohr}\)
using
\[
  \SI{1}{\bohr}=\SI{0.52917721067}{\angstrom}.
\]
For each element it reads a project-local D4 parameter row from
\texttt{PhysicalConstants.h}: electronegativity \(\chi\), hardness
\(\eta\) (named \texttt{gam} in the code), a coordination-number
electronegativity shift \(\kappa\), covalent radius \(r_{\mathrm{cov}}\), and
Gaussian charge width \(r_{\mathrm{g}}\). The values used for protein elements
are:
\begin{center}
\small
\begin{tabular}{crrrrr}
\hline
Elem. & \(\chi\) / \si{\hartree} & \(\eta\) / \si{\hartree}
& \(\kappa\) / \si{\hartree} & \(r_{\mathrm{cov}}\) / \si{\bohr}
& \(r_{\mathrm{g}}\) / \si{\bohr} \\
\hline
H & -0.35015861 & 0.47259288 & -0.19793756 & 0.80628 & 1.61478\\
C & -0.04726052 & 0.25364654 &  0.14216971 & 1.51718 & 2.49988\\
N &  0.11527249 & 0.28022740 &  0.15169154 & 1.42165 & 2.23456\\
O &  0.25136810 & 0.36515829 &  0.14510449 & 1.24854 & 1.89247\\
S &  0.10789083 & 0.25140725 &  0.15916035 & 2.00000 & 3.29733\\
\hline
\end{tabular}
\end{center}
An unrecognized element uses the fallback row
\((0,\;0.30000000,\;0,\;1.50000,\;2.50000)\) in the same units.

The first computed quantity is a smooth coordination number. For each atom pair
within the configured cutoff \(R_{\mathrm{cut}}=\SI{25.0}{\angstrom}\), the code
adds
\begin{equation}
  c_{ij} =
  \frac{1}{2}\operatorname{erfc}
  \left[
    k_{\mathrm{CN}}
    \left(
      \frac{R_{ij}}{r_{\mathrm{cov},i}+r_{\mathrm{cov},j}} - 1
    \right)
  \right],
  \qquad k_{\mathrm{CN}}=7.5 .
  \label{eq:cnpair}
\end{equation}
\(R_{ij}\) is the interatomic distance in \(\si{\bohr}\). The complementary
error function \(\operatorname{erfc}\) makes the count change smoothly from near
one for close bonded neighbours to near zero for separated atoms. The total
coordination number is
\begin{equation}
  \mathrm{CN}_i = \sum_{j\ne i} c_{ij}.
  \label{eq:cn}
\end{equation}
Pairs beyond \(R_{\mathrm{cut}}\) are skipped.

The coordination number shifts the electronegativity:
\begin{equation}
  \chi^{\mathrm{eff}}_i =
  \chi_i + \kappa_i\sqrt{\mathrm{CN}_i+10^{-14}} .
  \label{eq:chieff}
\end{equation}
The small \(10^{-14}\) term is a guard against evaluating the square root at
zero. \(\chi^{\mathrm{eff}}_i\) has units of \(\si{\hartree}\) per elementary
charge in atomic-unit bookkeeping.

The charge vector is then obtained by minimizing a quadratic energy subject to
a total-charge constraint:
\begin{equation}
  E(\bm q)=
  \sum_i \chi^{\mathrm{eff}}_i q_i
  + \frac{1}{2}\sum_i A_{ii}q_i^2
  + \sum_{i<j}A_{ij}q_iq_j ,
  \qquad \sum_i q_i = Q .
  \label{eq:energy}
\end{equation}
\(q_i\) is measured in elementary charges, and \(Q\) is the configured total
charge, defaulting to \(\SI{0}{\elementarycharge}\). The matrix diagonal is
\begin{equation}
  A_{ii}=\eta_i+\frac{\sqrt{2/\pi}}{r_{\mathrm{g},i}},
  \label{eq:adiag}
\end{equation}
where the second term is the self-Coulomb energy of the Gaussian charge
distribution used by the model. For \(i\ne j\), the implementation uses the
Ohno--Klopman form
\begin{equation}
  A_{ij}=\gamma_{ij}(R_{ij})
  =
  \frac{1}{\sqrt{R_{ij}^2+1/(\eta_i\eta_j)}} .
  \label{eq:ohno}
\end{equation}
This is the code's off-diagonal Coulomb kernel. It differs from the dftd4
reference implementation, which uses a Gaussian-charge
\(\operatorname{erf}(\alpha R)/R\) form. At long range
Eq.~\eqref{eq:ohno} approaches the atomic-unit Coulomb factor \(1/R_{ij}\). At
short range it remains finite, approaching \(\sqrt{\eta_i\eta_j}\) as
\(R_{ij}\to0\).

The equalization condition follows by differentiating Eq.~\eqref{eq:energy}.
Define the marginal charge energy
\begin{equation}
  \mu_i(\bm q)=\frac{\partial E}{\partial q_i}
  =
  \chi^{\mathrm{eff}}_i+\sum_j A_{ij}q_j .
  \label{eq:mu}
\end{equation}
Moving a small amount of charge from atom \(j\) to atom \(i\) changes the energy
to first order by \((\mu_i-\mu_j)\,dq\). At the constrained minimum those
exchange costs must match, so the code solves
\begin{equation}
  \begin{bmatrix}
    \bm A & \bm 1\\
    \bm 1^{\mathsf T} & 0
  \end{bmatrix}
  \begin{bmatrix}
    \bm q\\ \lambda
  \end{bmatrix}
  =
  \begin{bmatrix}
    -\bm\chi^{\mathrm{eff}}\\ Q
  \end{bmatrix}.
  \label{eq:kkt}
\end{equation}
\(\bm 1\) is the vector of ones, and the Lagrange multiplier \(\lambda\) is the
common offset that enforces the total charge. The augmented system is the
algebraic form of equal marginal energy plus charge conservation.

Rather than factor the bordered matrix in Eq.~\eqref{eq:kkt}, the code factors
\(\bm A\) with Eigen's Cholesky routine \texttt{LLT}. Cholesky writes a
symmetric positive-definite matrix as \(\bm L\bm L^{\mathsf T}\), turning later
linear solves into triangular substitutions. It then solves
\begin{equation}
  \bm A\bm u=\bm\chi^{\mathrm{eff}},
  \qquad
  \bm A\bm v=\bm 1,
  \label{eq:uv}
\end{equation}
and computes
\begin{equation}
  \lambda =
  -\frac{Q+\bm 1^{\mathsf T}\bm u}{\bm 1^{\mathsf T}\bm v},
  \qquad
  \bm q=-(\bm u+\lambda\bm v).
  \label{eq:solve}
\end{equation}
If the factorization fails, or if the denominator in Eq.~\eqref{eq:solve} is
smaller than \(10^{-30}\), the result is not attached.

After the solve, the code removes the floating-point charge residual by a
uniform shift:
\begin{equation}
  q_i \leftarrow q_i -
  \frac{\sum_j q_j-Q}{N}.
  \label{eq:residual}
\end{equation}
This preserves charge differences while making the pre-clamp sum equal to
\(Q\) to roundoff. A guard then clamps any stored atom charge with
\(|q_i|>\SI{2.0}{\elementarycharge}\) to
\(\pm\SI{2.0}{\elementarycharge}\). If that guard fires, the stored charge sum
can move away from \(Q\). The code records clamp events and charge statistics
for audit output.

\section{Inputs and output}

EEQ consumes atom elements, atom positions in \(\si{\angstrom}\), and four TOML
configuration values: \(Q=\SI{0}{\elementarycharge}\),
\(k_{\mathrm{CN}}=7.5\), \(R_{\mathrm{cut}}=\SI{25.0}{\angstrom}\), and the
\(\SI{2.0}{\elementarycharge}\) charge clamp. It has no declared
\texttt{ConformationResult} dependency and no external binary.

It does not consume MOPAC charges, AIMNet2 charges or embeddings, ff14SB
charges, APBS electrostatic fields, or DSSP secondary-structure labels. Those
data sources feed other parts of the pipeline: MOPAC and AIMNet2 provide
separate charge models, ff14SB charges feed the force-field Coulomb field, APBS
provides a solvated field and field gradient, and DSSP supplies secondary
structure labels. In the current code, no field calculator reads
\texttt{eeq\_charge}.

The feature output is scalar. \texttt{eeq\_charges.npy} has shape \((N,)\) and
stores \(q_i\) in elementary charges. \texttt{eeq\_cn.npy} has shape \((N,)\)
and stores \(\mathrm{CN}_i\). Across a trajectory,
\texttt{EeqWelfordTrajectoryResult} accumulates means, standard deviations,
minima, maxima, frame-to-frame changes, and charge-change rates for \(q_i\).
None of these values is a shielding. Calibration or later modelling is the step
that can test whether this geometry-responsive charge signal correlates with
DFT shielding changes.

\end{document}
